\documentclass[10pt,a4paper]{amsart}
\usepackage[margin=19mm]{geometry}
\usepackage{amsmath,amssymb,mathtools}
\usepackage[scr=rsfs,cal=euler]{mathalfa}
\usepackage{xfrac}
\usepackage[unicode,hidelinks,bookmarks=false]{hyperref}
\newcommand{\ie}{i.e.\ }
\newcommand{\cf}{cf.\ }
\newcommand{\resp}{respectively\ }
\newcommand{\Subsection}{\subsection}
\providecommand{\nobreakdash}{\nobreak\hbox{-}}
\setlength{\emergencystretch}{2em}
\multlinegap0pt
\usepackage{amssymb}
\usepackage{float}
\usepackage[shortlabels]{enumitem}
\usepackage{xcolor}
\newtheorem{theorem}{Theorem}
\newtheorem{corollary}{Corollary}
\newtheorem{proposition}{Proposition}
\newcommand{\colo}[2]{(1^{#1}, 2^{#2})}
\newcommand{\ex}[3]{\operatorname{ex}^{#1}_{#2}\hspace{-3pt}{#3}}
\title{Between proper and square coloring of planar graphs, hardness and extremal graphs}
\author{Thomas Del\'epine}
\date{Source: arXiv:2602.13037v1 (CC BY 4.0)}
\begin{document}
\maketitle

\noindent\emph{Source and licence.} Between proper and square coloring of planar graphs, hardness and extremal graphs. Version arXiv:2602.13037v1 (CC BY 4.0), distributed by arXiv under the Creative Commons Attribution 4.0 International licence, http://creativecommons.org/licenses/by/4.0/, as recorded in the arXiv licence metadata for this exact version and shown on the abstract page https://arxiv.org/abs/2602.13037v1. Full licence notice: \texttt{licenses/arxiv-2602.13037-CC-BY-4.0.txt}.\par\smallskip

\noindent\textit{Editorial scope.} Counted unit: the article's proposition prop:upperbound\_colo\_for\_forests\_and\_triangle-free\_outerplanar\_graphs (source0.tex lines 512--514). The article's complete original proof of it is delivered with it (source0.tex lines 516--533, one \texttt{proof} environment of 18 lines). No mathematics is written, repaired or re-derived: the statement and the proof are the article's own lines, in the article's own notation, and no retained line is substituted or re-flowed.  Retained uncounted as the source-local context the proof needs: lines 444--446; lines 500--502.  Nothing was omitted from any retained span, so the package carries no omission record.  No retained line contains a citation, so the wrapper carries no bibliography and no bibliography entry.  No retained line contains a figure environment, an image-inclusion command or drawing code, so no image is delivered and the package declares no asset dependency.  Editorial formatting only, in addition: an AMS \texttt{amsart} preamble that restates the article's own declarations and macros used by the retained lines, namely the environments \texttt{theorem}, \texttt{corollary}, \texttt{proposition} on separate counters, together with the standard packages those lines need.  The retained environments print in this document as Theorem 1, Corollary 1, Proposition 1 in order of appearance, whereas the article prints the same items with the numbering of its own full document.  The complete native source of the article as distributed in the arXiv e-print for this exact version is bundled unchanged as \texttt{sources/194605/source0.tex}.
\begin{theorem}\label{thm:coloring_k_degenerate_graphs}
   Let $\mathcal{D}_k$ be the class of $k$-degenerate graphs. Then $\ex{k}{\mathcal{D}_k}{(n)} \leq 4k\sqrt{k + 1}\sqrt{n}$.
\end{theorem}

\begin{corollary}\label{coro:vertex_cover_coro}
    Let $G$ be an outer-planar graph of girth $g$ and of order $n$, then $G$ has a vertex cover of order at most $n(k + 1)/(2k + 1)$ where $k$ is such that $g = 2k$ or $g = 2k + 1$.
\end{corollary}

\begin{proposition}\label{prop:upperbound_colo_for_forests_and_triangle-free_outerplanar_graphs}
    Forests are $\colo{1}{4\sqrt{2}\sqrt{n}}$-colorable and triangle-free outerplanar graphs are $\colo{1}{4\sqrt{34/5}\sqrt{n} - 1}$-colorable.
\end{proposition}

\begin{proof}
    Forests are $1$-degenerate so we can apply Theorem~\ref{thm:coloring_k_degenerate_graphs} to $\colo{1}{4\sqrt{2}\sqrt{n}}$-color every forests. 
    Let $G$ be a triangle-free outerplanar graph, and let $S$ be the set of vertices of degree at least $\alpha \sqrt{n}$ for some real number $\alpha$ to be optimized later. Since $G$ is outerplanar, it is $2$-degenerate, so $G$ has at most $2n$ edges and therefore $|S| \leq 4\sqrt{n}/\alpha$. 

    In order to color $G$, our first goal is to construct a set $S' \subseteq N_G[S]$ such that $S \subseteq S'$, $|S'| = \Theta(|S|)$ and $N_G[S] - S'$ is an independent set. If $|S| = 1$, then since $G$ is triangle-free, $N_G[S] - S'$ is an independent set. Otherwise, we claim that if $|S| \geq 2$, then $|N_G[S]| \le 5|S|-4$.
    Indeed, we proceed by induction on $\kappa_2$, the number of $2$-connected components of $G[N_G[S]]$. If $\kappa_2 = 1$, then the outerface of $G[N_G[S]]$ is a cycle and since $S$ is a dominating set of $G[N_G[S]]$ and $|S| \ge 2$, we have $|N_G[S]| \le 3|S| \le  5|S| - 4$. Now assume that $\kappa_2 \geq 2$ and let $B$ be a leaf block of $G[N_G[S]]$, that is a $2$-connected component of $G[N_G[S]]$ with exactly one cut-vertex $v$. Let $m = |S \cap B|$, we then have, 
        \[|G[N_G[S]]| = |G[N_G[S] - (B - v)]| + |G[B]| - 1\]
    since $G[N_G[S] - (B - v)]$ and $G[B]$ only interest on $v$. So by induction hypothesis, 
        \[|G[N_G[S]]| \le 5(|S| - (m - 1)) - 4 + 5m - 4 - 1 = 5|S| - 4\]
    as desired. Now consider the set 
    \[V' = \{v | \exists svv's' \textrm{ a path of } G \textrm{ with } s, s' \in S \textrm{ and } v, v' \in N_G[S] - S\}\]
    Then by the previous claim, and since $V' \subseteq N_G[S] - S$, $|V'| \le 4|S| - 4 \le 4|S|$ so, in particular, there exists a set $C \subseteq V'$ that is a vertex cover of $G[V']$ of size at most $12|S|/5$ since by Corollary \ref{coro:vertex_cover_coro}, every outerplanar graph or order $n$ has a vertex cover of size at most $3n/5$. We thus define $S'$ to be $S \cup C$ so $|S'| \le 17|S|/5$. Assume that $N_G[S] - S'$ is not an independent set. Then there exists $u$, $v \in N_G[S] - S'$ such that $uv$ is an edge of $G[N_G[S]]$. Since $S$ is a dominating set of $N_G[S]$ and $G$ is triangle-free, there must exist $s$, $s' \in S$ such that $suvs'$ is a path in $G[N_G[S]]$. But $u$ and $v$ belong to $V'$ and $uv$ is not covered by $C$ and this is a contradiction.

    We can now color $G$. We first construct the $2$-tree $T$ that is a supergraph of $G$ and that we will color. Since $G$ is outerplanar, it has treewidth at most $2$ so $T$ exists. We first color every vertex from $S'$ with a unique distance-2 color each, at the cost of at most $|S'| \leq 17|S|/5 \le 68\sqrt{n}/(5\alpha)$ distance-2 colors. Then, we color every vertex in $N_G(S) - S'$ with a distance-1 color. This pre-coloring of $T$ might not be valid, but it is valid when we reduce it to $G$ since as previously observed, $N_G(S) - S'$ is an independent set in $G$ and each distance-2 color is used once. Then, we color the remaining vertices of $T$ following a $2$-degenerate ordering provided by the construction of this $2$-tree. 
    Let $v_1, \dots, v_n$ be an ordering of the vertices and $T_i$ be the subgraph of $T$ induced by $\{v_1, \dots, v_i\}$ such that for every $i$, $\deg_{T_i}(v_i) = 2$. By definition of a $2$-tree, the two neighbors of $v_i$ are neighbors themselves. Let $A$ be a set of distance-2 colors distinct from the distance-2 colors used to color vertices in $S'$ whose size will be determined later. If $v_1 \not\in N_G[S]$ was precolored, then we do not recolor it. Otherwise, we can color it with any distance-2 color from $A$. 
    Then, assume that $i \ge 2$ and that we managed to color $T_{i - 1}$ while respecting the precoloring. We thus need to color $v_i$ in order to extend our coloring up to $T_i$.
    As previously observed, the two neighbors of $v_i$ in $G_i$ are neighbors themselves so they do not have the same color. Therefore, by adding $v_i$ to $T_{i - 1}$, we do not create conflicts. If $v_i$ was precolored, then we do not recolor it. Otherwise, the neighbors of $v_i$ are of degree at most $\alpha \sqrt{n} - 1$ in $G_{i - 1}$ and so at most $2\alpha \sqrt{n} - 2$ distance-2 colors are unavailable for $v_i$. Therefore, if $|A| \geq 2\alpha \sqrt{n} - 1$, we can color $T$ such that there is no conflict involving the distance-2 colors. When reducing this coloring to $G$, there is still no conflict involving distance-2 colors since $G$ is a subgraph of $T$ and there is no conflict involving distance-1 colors since the vertices colored with the distance-1 color induce an independent set. In total, we used at most $|S'| + |A| \le 68\sqrt{n}/(5\alpha) +  2\alpha \sqrt{n} - 1$ distance-2 colors and this quantity is minimized for $\alpha = \sqrt{34/5}$ in which case we managed to $\colo{1}{4\sqrt{34/5}\sqrt{n} - 1}$-color $G$ as desired.
\end{proof}

\end{document}
